\textbf{SAM and variants.} SAM \cite{foret2020sharpness} minimises $\max_{\|\epsilon\|\le\rho}L(w+\epsilon)$ with a one-step approximation of the inner maximum. ASAM \cite{kwon2021asam} makes the neighbourhood adaptive to parameter scale, addressing the dependence of fixed-radius sharpness on re-parameterisation, and ESAM \cite{du2021efficient} reduces SAM's cost of two gradient evaluations per step. We reimplement only the original SAM; the scale dependence that ASAM addresses matters for our sharpness measure (Sec.~\ref{sec:limitations}).

\textbf{Understanding SAM.} Andriushchenko and Flammarion \cite{andriushchenko2022understanding} analyse the implicit bias of SAM (diagonal linear networks, plus convergence results) and argue that PAC-Bayes and flat-minima explanations are incomplete. Wen et al.\ \cite{wen2022how} pin down which notion of sharpness SAM actually regularises. For label noise specifically, Baek et al.\ \cite{baek2024why} report that SAM's largest gains occur under label noise, that the peak occurs with early stopping rather than at convergence, and decompose the robustness into a logit effect (up-weighting clean examples) and a Jacobian effect. We use no early stopping, so our results concern the end of training only.

\textbf{Sharpness and generalisation.} Keskar et al.\ \cite{keskar2016large} associated large-batch training with sharp minima and a larger generalisation gap; Jiang et al.\ \cite{jiang2019fantastic} compared over 40 complexity measures across thousands of networks; Andriushchenko et al.\ \cite{andriushchenko2023modern} report that (adaptive) sharpness does not correlate well with generalisation across modern settings, but rather with training parameters such as the learning rate. Zhang et al.\ \cite{zhang2016understanding} showed that standard networks fit random labels, which motivates our measurement of how many flipped labels each optimiser memorises.
